\section{一般化線形モデル}
\subsection{授業内容}
\subsubsection{科目の中でのこのコマの位置づけ}
\subsubsection{コマ主題細目}
ここまで要因計画が線形モデルと同一であること，すなわち一般線形モデルについて議論されてきた。あらためて回帰分析の確率モデルを考えると，平均に構造を入れたモデルという意味で同じであることが確認できる。ここで確率分布を違う形に変えることでより一般的な線形モデル，一般化線形モデルに拡張することができる。確率モデルによっては結果変数の型の違いによって確率分布が変わり，確率分布によってはパラメータの取りうる範囲が定まるのでリンク関数によって変換する必要があること，結果を解釈するときはリンク関数を経由して分析されていることなどに注意が必要である。まずはパラメータの数が少ないロジスティック回帰，ポアソン回帰について学ぶ。
\begin{description}
    \item[回帰モデルの確率モデル] 正規分布の平均構造を導入するという意味で，回帰分析のベイズ推定はこれまでと同様に実施，解釈することができる。事後分布や事後予測分布などを使って，最尤推定のモデルと異なる点を確認しておく。
    \item[データに合わせた確率分布] データの形によっては確率分布の形を変える必要がある。まずはベルヌーイ分布によるロジスティック回帰分析を通じて，確率分布関数を選択できること，そのためにパラメータの形をリンク関数を経由し得て変換することを学ぶ。
    \item[リンク関数とパラメータの解釈] リンク関数を介して線形モデルを考えるので，独立変数が一単位増えることがそのまま従属変数が一単位増えることにはつながらない。これらの関係を知るために，事後予測分布などを出しながら考えていく必要がある。
    \item[ポアソン分布による回帰] 応用例として，ポアソン分布による回帰分析を考える。リンク関数や解釈の際の逆リンク関数の当てはめ方について学ぶ。

\end{description}
\subsubsection{キーワード}
\begin{itemize}
    \item 一般化線形モデル
    \item ベルヌーイ分布
    \item ポアソン分布
    \item ロジスティック回帰分析
    \item リンク関数
    \item ポアソン回帰分析
\end{itemize}
\subsection{授業情報}
\paragraph{コマの展開方法}
講義/遠隔/演習
\subsubsection{予習・復習課題}
\paragraph{予習}確率変数が連続的か，離散的かということがどういう違いであるのかについて，データ解析基礎の確率に関する資料などを参考に確認しておく。また尺度水準による数値データの分類についても見直しておくと良い。
\paragraph{復習}離散的なデータについて，身近な例を考えてみると良い。また今回導入した分布関数以外にも応用を考えることができるから，確率分布とリンク関数の一覧を参考にさまざまなモデルに思いを馳せてみると良い。